\section{Applications to definite integrals}

For rational functions of $\sin\theta$ and $\cos\theta$, the substitution $z=e^{\mathrm{i}\theta}$ maps $[0,2\pi]$ to the unit circle and gives
\begin{equation}
\cos\theta=\frac{z+z^{-1}}{2},\qquad
\sin\theta=\frac{z-z^{-1}}{2\mathrm{i}},\qquad d\theta=\frac{dz}{\mathrm{i}z}.
\end{equation}
For example, the pole inside the unit circle yields
\begin{equation}
\int_0^{2\pi}\frac{d\theta}{1+a\cos\theta}=\frac{2\pi}{\sqrt{1-a^2}},\qquad |a|<1.
\end{equation}
Indeed, after substitution the poles are
\begin{equation}
z_{\pm}=-\frac{1\pm\sqrt{1-a^2}}{a}.
\end{equation}
For $|a|<1$, exactly one lies inside $|z|=1$. Its simple residue produces the displayed integral. This unit-circle method is the standard treatment of every rational trigonometric integral of this form.

For integrals on the whole real line, a large semicircle in the upper half-plane is natural when the arc contribution vanishes. The residue theorem then gives
\begin{equation}
\int_{-\infty}^{\infty}f(x)\,dx=2\pi\mathrm{i}\sum_{\operatorname{Im}z_j>0}\operatorname{Res}(f;z_j).
\end{equation}

\begin{figure}[htb]
    \centering
    \scalebox{0.72}{\subimport{tikz/}{semicircle.tex}}
    \caption{Upper-half-plane semicircle used for real-axis integrals.}
    \label{fig:ch06-semicircle}
\end{figure}

Taking $f(z)=1/(1+z^2)$ leaves only the pole at $z=\mathrm{i}$ in the upper half-plane and gives $\int_{-\infty}^{\infty}dx/(1+x^2)=\pi$. With an exponential factor, Jordan's lemma similarly gives, for $a,m>0$,
\begin{equation}
\int_0^\infty\frac{\cos(mx)}{x^2+a^2}\,dx=\frac{\pi}{2a}e^{-am}.
\end{equation}

For the exponential contour method, assume $f$ is analytic apart from finitely many upper-half-plane poles and decreases sufficiently fast. On $z=Re^{\mathrm{i}\theta}$,
\begin{equation}
|e^{\mathrm{i}mz}|=e^{-mR\sin\theta}.
\end{equation}
Using $\sin\theta\geq2\theta/\pi$ on $[0,\pi/2]$ proves that the arc integral tends to zero. This is Jordan's lemma and gives
\begin{equation}
\int_{-\infty}^{\infty}f(x)e^{\mathrm{i}mx}\,dx
=2\pi\mathrm{i}\sum_{\operatorname{Im}z_j>0}\operatorname{Res}\left(e^{\mathrm{i}mz}f(z);z_j\right).
\end{equation}
Even cosine and odd sine integrals reduce to this formula by symmetry.

For the arc estimate, suppose $|f(z)|<\varepsilon$ on the large semicircle. Then
\begin{align}
\left|\int_{C_R}f(z)e^{\mathrm{i}mz}\,dz\right|
&\leq2\varepsilon R\int_0^{\pi/2}e^{-mR\sin\theta}\,d\theta\nonumber\\
&\leq2\varepsilon R\int_0^{\pi/2}e^{-2mR\theta/\pi}\,d\theta
<\frac{\pi\varepsilon}{m}.
\end{align}
This proves Jordan's lemma. If $F$ is even, then
\begin{equation}
\int_0^\infty F(x)\cos(mx)\,dx
=\frac12\int_{-\infty}^{\infty}F(x)e^{\mathrm{i}mx}\,dx;
\end{equation}
the analogous statement holds for odd sine integrals.

Parameter differentiation provides a complementary method. From
\begin{equation}
\int_0^\infty e^{-ax}\cos(bx)\,dx=\frac{a}{a^2+b^2},
\end{equation}
differentiation with respect to $a$ gives
\begin{equation}
\int_0^\infty xe^{-ax}\cos(bx)\,dx=\frac{a^2-b^2}{(a^2+b^2)^2}.
\end{equation}
Adding a damping factor to $\int_0^\infty\sin x/x\,dx$ and removing it at the end gives the classic result $\pi/2$. A keyhole contour gives the beta-integral form
\begin{equation}
\int_0^\infty\frac{x^{\alpha-1}}{1+x}\,dx=\frac{\pi}{\sin(\pi\alpha)},\qquad 0<\alpha<1.
\end{equation}